% !TEX program = pdflatex
% =====================================================================
%  Agrégation interne de mathématiques -- Session 2026
%  Première épreuve écrite (EAI MAT 1)
%  Corrigé rédigé : Vrai/Faux, Partie I, Partie II
%  + encadrés "Erreurs fréquentes à éviter"
% =====================================================================
\documentclass[11pt,a4paper]{article}

\usepackage[utf8]{inputenc}
\usepackage{amsmath,amssymb,amsthm,mathtools}
\usepackage{stmaryrd}                 % \llbracket, \rrbracket
\usepackage[margin=2.3cm]{geometry}
\usepackage{xcolor}
\usepackage{enumitem}
\usepackage{titlesec}
\usepackage[most]{tcolorbox}
\usepackage{hyperref}

% ------------------------- Palette -----------------------------------
\definecolor{bleuP}{RGB}{25,54,102}       % bleu profond (titres, énoncés)
\definecolor{bleuClair}{RGB}{240,244,250} % fond des énoncés
\definecolor{vertP}{RGB}{22,101,52}       % vert (corrigés)
\definecolor{vertClair}{RGB}{243,250,245}
\definecolor{orangeP}{RGB}{178,74,10}     % orange (pièges)
\definecolor{orangeClair}{RGB}{255,247,237}
\definecolor{grisP}{RGB}{90,90,90}

\hypersetup{colorlinks=true, linkcolor=bleuP, urlcolor=bleuP}

% ------------------------- Titres ------------------------------------
\titleformat{\section}
  {\Large\bfseries\color{bleuP}}{\thesection.}{0.6em}{}
  [\vspace{2pt}{\color{bleuP}\titlerule[1.2pt]}]
\titleformat{\subsection}
  {\large\bfseries\color{bleuP}}{\thesubsection.}{0.6em}{}

% ------------------------- Boîtes ------------------------------------
\newtcolorbox{enonce}[1]{%
  enhanced, breakable,
  colback=bleuClair, colframe=bleuP,
  boxrule=0.8pt, arc=2mm,
  left=8pt, right=8pt, top=6pt, bottom=6pt,
  title={\bfseries #1},
  fonttitle=\color{white},
  colbacktitle=bleuP,
  attach boxed title to top left={yshift=-2mm, xshift=4mm},
  boxed title style={arc=1.5mm, boxrule=0pt}}

\newtcolorbox{pieges}{%
  enhanced, breakable,
  colback=orangeClair, colframe=orangeP,
  boxrule=0.8pt, arc=2mm,
  left=8pt, right=8pt, top=6pt, bottom=6pt,
  title={\bfseries Erreurs fréquentes à éviter},
  fonttitle=\color{white},
  colbacktitle=orangeP,
  attach boxed title to top left={yshift=-2mm, xshift=4mm},
  boxed title style={arc=1.5mm, boxrule=0pt}}

\newtcolorbox{conseil}{%
  enhanced, breakable,
  colback=vertClair, colframe=vertP,
  boxrule=0.8pt, arc=2mm,
  left=8pt, right=8pt, top=6pt, bottom=6pt,
  title={\bfseries Conseils de rédaction pour l'épreuve},
  fonttitle=\color{white},
  colbacktitle=vertP,
  attach boxed title to top left={yshift=-2mm, xshift=4mm},
  boxed title style={arc=1.5mm, boxrule=0pt}}

% Corrigé : filet vertical vert à gauche
\newenvironment{corrige}
  {\par\medskip\noindent
   \begin{tcolorbox}[blanker, breakable,
     borderline west={2.5pt}{0pt}{vertP},
     left=10pt, top=2pt, bottom=2pt]
   {\color{vertP}\bfseries Corrigé.}\hspace{0.5em}\ignorespaces}
  {\hfill$\blacksquare$\end{tcolorbox}\medskip}

% ------------------------- Macros ------------------------------------
\newcommand{\og}{``}
\newcommand{\fg}{''}
\newcommand{\K}{\mathbb{K}}
\newcommand{\Z}{\mathbb{Z}}
\newcommand{\Q}{\mathbb{Q}}
\newcommand{\R}{\mathbb{R}}
\newcommand{\N}{\mathbb{N}}
\newcommand{\Fp}{\mathbb{F}_p}
\newcommand{\Mn}[1][n]{\mathcal{M}_{#1}}
\newcommand{\GLn}[1][n]{\mathrm{GL}_{#1}}
\newcommand{\On}[1][n]{\mathrm{O}_{#1}}
\newcommand{\Sym}[1][n]{S_{#1}}          % matrices symétriques S_n(Z)
\newcommand{\Perm}[1][n]{\mathcal{S}_{#1}} % permutations
\newcommand{\ii}[2]{\llbracket #1,#2\rrbracket}
\newcommand{\dg}[1]{\mathrm{d}^{\circ}\!\left(#1\right)}
\newcommand{\T}{\mathsf{T}}               % transposée : M^\T
\DeclareMathOperator{\Ker}{Ker}
\DeclareMathOperator{\Ima}{Im}
\DeclareMathOperator{\rg}{rang}
\DeclareMathOperator{\vect}{vect}
\DeclareMathOperator{\Mat}{Mat}
\DeclareMathOperator{\Diag}{Diag}
\newcommand{\red}[1]{\overline{#1}}       % réduction mod p

\setlist[enumerate,1]{label=(\alph*), leftmargin=2em}

\title{\vspace{-1.2cm}{\color{bleuP}\bfseries\LARGE
Agrégation interne de mathématiques --- Session 2026\\[4pt]
\Large Première épreuve écrite (EAI MAT 1)}\\[8pt]
{\large Corrigé rédigé : Vrai/Faux, Partie I et Partie II,\\
avec les erreurs fréquentes à éviter}}
\author{}
\date{}

\begin{document}
\maketitle
\vspace{-1.4cm}

\begin{center}
\begin{minipage}{0.92\textwidth}\small\itshape\color{grisP}
Le problème porte sur l'action de conjugaison de $\On(\Q)$ sur
$\Sym(\R)\cap\Mn(\Z)$. Le Vrai/Faux et les parties I et II, corrigés
ci-dessous, mettent en place les outils : réduction modulo $p$,
groupe $\GLn(\Z)$, matrices de permutation et de transvection,
puis le résultant $\Delta(R)=\det(\psi_R)$ et l'étude fine de
$\Ker(\varphi_{R,S})$. Chaque question est suivie, lorsque c'est
pertinent, d'un encadré signalant les pièges classiques relevés
dans les rapports de jury.
\end{minipage}
\end{center}

\begin{conseil}
\begin{itemize}[leftmargin=1.4em, itemsep=2pt]
\item \textbf{Le Vrai/Faux se justifie.} L'énoncé le dit explicitement :
une réponse non argumentée (ou un contre-exemple non vérifié) ne
rapporte rien. Pour un \emph{faux}, un contre-exemple explicite et
\emph{vérifié} ; pour un \emph{vrai}, une démonstration complète.
\item \textbf{Citer les questions utilisées.} Le sujet est très
structuré ; écrire \og d'après la question 6(b)\fg{} est attendu et
valorisé. \emph{Remarque :} on écrira ici simplement
\emph{d'après 6(b)}.
\item \textbf{Attention à la caractéristique.} Le corps $\K$ est
quelconque : plusieurs questions (VF4, 17, 18) sont précisément
conçues pour piéger les réflexes de la caractéristique nulle
($\dg{R'}=\dg{R}-1$, \og $k\neq 0$ \fg, etc.).
\item \textbf{\og Soigneusement \fg{} signifie coefficient par
coefficient.} Aux questions 9 et 10, un dessin ou un exemple
$3\times 3$ ne suffit pas : on calcule le coefficient générique du
produit avec le symbole de Kronecker.
\item \textbf{Vérifier les appartenances.} Quand on exhibe un élément
d'un noyau ou d'un espace $\K_{d-1}[X]$, il faut contrôler les
\emph{degrés} : c'est le point le plus souvent omis en question 15.
\end{itemize}
\end{conseil}

% =====================================================================
\section{Vrai ou faux ?}
% =====================================================================

% ------------------------------------------------ VF 1
\begin{enonce}{Question 1}
Soient $p$ un nombre premier et $A\in\Mn(\Z)$.
Affirmation : \emph{dans $\Fp$, on a
$\red{\det(A)}=\det\big(\red{A}\big)$.}
\end{enonce}

\begin{corrige}
\textbf{L'affirmation est vraie.}
Notons $\pi:\Z\to\Fp$, $a\mapsto\red{a}$, la réduction modulo $p$ :
c'est un \emph{morphisme d'anneaux} (surjectif). Écrivons la formule de
Leibniz pour $A=[a_{i,j}]$ :
\[
\det(A)\;=\;\sum_{\sigma\in\Perm}\varepsilon(\sigma)
\prod_{i=1}^{n}a_{\sigma(i),i}.
\]
Comme $\pi$ est additif et multiplicatif et vérifie $\pi(-x)=-\pi(x)$,
il commute avec les sommes, les produits et les signes ; en appliquant
$\pi$ membre à membre :
\[
\red{\det(A)}
=\sum_{\sigma\in\Perm}\varepsilon(\sigma)
\prod_{i=1}^{n}\red{a_{\sigma(i),i}}
=\det\big([\,\red{a_{i,j}}\,]\big)
=\det\big(\red{A}\big),
\]
où l'on a reconnu la formule de Leibniz du déterminant dans le corps
$\Fp$ (le signe $\varepsilon(\sigma)=\pm 1$ est envoyé sur
$\pm\red{1}$).
\end{corrige}

\begin{pieges}
\begin{itemize}[leftmargin=1.4em, itemsep=2pt]
\item Invoquer la \og linéarité \fg{} de la réduction modulo $p$ : le
déterminant n'est pas linéaire en $A$. C'est le caractère de
\emph{morphisme d'anneaux} (compatibilité aux sommes \emph{et} aux
produits) qui fait fonctionner l'argument, via une formule explicite
du déterminant.
\item Raisonner avec $\det(A)=\prod\lambda_i$ (valeurs propres) : les
valeurs propres n'ont aucune raison d'être entières ni de se réduire
convenablement.
\item Oublier de dire ce que devient $\varepsilon(\sigma)$ dans $\Fp$.
\end{itemize}
\end{pieges}

% ------------------------------------------------ VF 2
\begin{enonce}{Question 2}
Soient $E$ et $F$ deux espaces vectoriels de bases respectives
$(e_i)_{1\le i\le m}$ et $(f_j)_{1\le j\le n}$.
Affirmation : \emph{$(e_i,f_j)_{\substack{1\le i\le m\\ 1\le j\le n}}$
est une base de $E\times F$.}
\end{enonce}

\begin{corrige}
\textbf{L'affirmation est fausse.}
La famille proposée est constituée des $mn$ couples
$(e_i,f_j)\in E\times F$. Montrons qu'elle n'engendre jamais
$E\times F$ : le vecteur $(e_1,0)$ n'appartient pas à son sous-espace
engendré. En effet, si
\[
(e_1,0)\;=\;\sum_{i=1}^{m}\sum_{j=1}^{n}\lambda_{i,j}\,(e_i,f_j)
\;=\;\Big(\sum_{i}\big(\textstyle\sum_{j}\lambda_{i,j}\big)e_i\,,\;
\sum_{j}\big(\textstyle\sum_{i}\lambda_{i,j}\big)f_j\Big),
\]
alors, par liberté de $(e_i)$ et de $(f_j)$ :
\[
\sum_{j}\lambda_{1,j}=1,\qquad
\sum_{j}\lambda_{i,j}=0\ (i\ge 2),\qquad
\sum_{i}\lambda_{i,j}=0\ (1\le j\le n).
\]
En sommant toutes les égalités de gauche, la somme totale
$\Sigma=\sum_{i,j}\lambda_{i,j}$ vaut $1$ ; en sommant celles de
droite, $\Sigma=0$ : contradiction. La famille n'est donc pas
génératrice, ce n'est pas une base.

\emph{Complément.} La bonne base de $E\times F$ est la famille de
cardinal $m+n$
\[
\big((e_1,0),\dots,(e_m,0),(0,f_1),\dots,(0,f_n)\big),
\]
et $\dim(E\times F)=m+n$. On notera d'ailleurs que pour $m,n\ge 2$
la famille de l'énoncé n'est même pas libre :
$(e_1,f_1)+(e_2,f_2)-(e_1,f_2)-(e_2,f_1)=(0,0)$.
\end{corrige}

\begin{pieges}
\begin{itemize}[leftmargin=1.4em, itemsep=2pt]
\item Se contenter de l'argument de cardinal $mn\neq m+n$ : il est en
défaut pour $m=n=2$ (où $mn=m+n=4$) ! Il faut un argument valable
pour \emph{tous} $m,n$, comme celui donné ci-dessus, ou traiter ce
cas à part.
\item Confondre avec le \emph{produit tensoriel} $E\otimes F$, dont
$(e_i\otimes f_j)$ est effectivement une base : c'est l'origine du
piège.
\item Répondre \og faux \fg{} sans justification, ou avec un
contre-exemple non vérifié.
\end{itemize}
\end{pieges}

% ------------------------------------------------ VF 3
\begin{enonce}{Question 3}
Soient $E$ de dimension finie, $f\in\mathcal{L}(E)$ diagonalisable et
$F$ un sous-espace vectoriel de $E$.
Affirmation : \emph{si $F$ est stable par $f$, alors $f$ induit un
endomorphisme diagonalisable de $F$.}
\end{enonce}

\begin{corrige}
\textbf{L'affirmation est vraie.}
Notons $g:F\to F$, $x\mapsto f(x)$, l'endomorphisme induit (bien
défini car $f(F)\subset F$).

Puisque $f$ est diagonalisable, il existe un polynôme
$P\in\K[X]$ \emph{scindé à racines simples} annulant $f$ (par exemple
$P=\prod_{\lambda\in\mathrm{Sp}(f)}(X-\lambda)$).

Montrons que $P(g)=0$. Une récurrence immédiate donne, pour tout
$k\in\N$ et tout $x\in F$, $g^{k}(x)=f^{k}(x)$ : c'est vrai pour
$k=0$, et si $g^k(x)=f^k(x)$ alors, comme $f^k(x)\in F$ (stabilité),
$g^{k+1}(x)=g\big(f^k(x)\big)=f^{k+1}(x)$. Par linéarité, pour tout
$x\in F$ :
\[
P(g)(x)=P(f)(x)=0 .
\]
Ainsi $g$ est annulé par un polynôme scindé à racines simples : par le
critère de diagonalisabilité (via le lemme des noyaux), $g$ est
diagonalisable.
\end{corrige}

\begin{pieges}
\begin{itemize}[leftmargin=1.4em, itemsep=2pt]
\item Écrire directement $P(g)=P(f)_{|F}$ sans justification : la
petite récurrence $g^k(x)=f^k(x)$ (qui utilise la stabilité à chaque
étape) est attendue.
\item Tenter la voie \og $F=\bigoplus (F\cap E_\lambda)$ \fg{} : le
résultat est vrai mais c'est précisément ce qu'il faut démontrer ;
beaucoup de copies l'affirment sans preuve. La voie du polynôme
annulateur est la plus économique.
\item Énoncer le critère de diagonalisabilité de façon incomplète
(oublier \og à racines simples \fg{} ou \og scindé \fg).
\end{itemize}
\end{pieges}

% ------------------------------------------------ VF 4
\begin{enonce}{Question 4}
Soient $\K$ un corps, $R\in\K[X]$ et $R'$ son polynôme dérivé.
Affirmation : \emph{si $R'=0$, alors $R$ est constant.}
\end{enonce}

\begin{corrige}
\textbf{L'affirmation est fausse} (pour un corps $\K$ quelconque).
Prenons $\K=\Fp$ avec $p$ premier et $R=X^{p}$. Alors
\[
R'=\red{p}\,X^{p-1}=\red{0}\cdot X^{p-1}=0,
\]
puisque $p\equiv 0\ [p]$, alors que $R$ n'est pas constant
($\dg{R}=p\ge 2$).

\emph{Complément.} L'affirmation est vraie si $\K$ est de
caractéristique nulle : si $R=\sum_{k=0}^{n}r_kX^k$ avec $R'=0$,
alors $k\,r_k=0$ pour tout $k\ge 1$, et comme $k\neq 0$ dans $\K$
on obtient $r_k=0$ pour $k\ge 1$. En caractéristique $p$, la
conclusion correcte est seulement : $R'=0$ si et seulement si $R$ est
un polynôme en $X^{p}$ (c'est l'objet de la question 34 du sujet).
\end{corrige}

\begin{pieges}
\begin{itemize}[leftmargin=1.4em, itemsep=2pt]
\item Réflexe \og $\dg{R'}=\dg{R}-1$ \fg{} : cette formule est fausse
en caractéristique $p$ (le coefficient dominant $n\,r_n$ peut
s'annuler).
\item Oublier que $\K$ est un corps \emph{quelconque} : l'énoncé ne dit
pas $\K=\R$ ou $\mathbb{C}$. Ce Vrai/Faux annonce la partie IV.
\item Donner le contre-exemple $X^p$ sans vérifier explicitement
$R'=0$ dans $\Fp[X]$.
\end{itemize}
\end{pieges}

% =====================================================================
\section{Partie I. Quelques résultats sur les matrices à coefficients
dans $\Z$}
% =====================================================================
Dans toute cette partie, $n\in\N^{*}$.

% ------------------------------------------------ Q5
\begin{enonce}{Question 5}
Soient $Q\in\On(\Q)$ et $A\in\Mn(\Z)$ telles que $Q^{\T}AQ\in\Mn(\Z)$.
On note $\widehat{Q}=\ell(Q)\,Q$ et on suppose que $p$ est un nombre
premier divisant $\ell(Q)$. Démontrer que
$\widehat{Q}^{\T}\widehat{Q}\equiv 0\,[p^{k}]$ et
$\widehat{Q}^{\T}A\widehat{Q}\equiv 0\,[p^{k}]$ pour $k\in\{1,2\}$.
\end{enonce}

\begin{corrige}
Posons $\ell=\ell(Q)$ ; par définition
$\widehat{Q}=\ell Q\in\Mn(\Z)$. Comme $Q\in\On(\Q)$, on a
$Q^{\T}Q=I_n$, d'où :
\[
\widehat{Q}^{\T}\widehat{Q}
=(\ell Q)^{\T}(\ell Q)
=\ell^{2}\,Q^{\T}Q
=\ell^{2}I_n .
\]
De même, par bilinéarité du produit matriciel :
\[
\widehat{Q}^{\T}A\widehat{Q}=\ell^{2}\,\big(Q^{\T}AQ\big),
\qquad\text{avec }Q^{\T}AQ\in\Mn(\Z)\text{ par hypothèse.}
\]
Ainsi, tout coefficient de $\widehat{Q}^{\T}\widehat{Q}$
(respectivement de $\widehat{Q}^{\T}A\widehat{Q}$) est de la forme
$\ell^{2}m$ avec $m\in\Z$. Or $p\mid\ell$, donc $p^{2}\mid\ell^{2}$,
et par conséquent $p^{2}\mid\ell^{2}m$ : chaque coefficient est
divisible par $p^{2}$, c'est-à-dire
\[
\widehat{Q}^{\T}\widehat{Q}\equiv 0\,[p^{2}]
\quad\text{et}\quad
\widehat{Q}^{\T}A\widehat{Q}\equiv 0\,[p^{2}].
\]
Enfin, la divisibilité par $p^{2}$ entraîne la divisibilité par $p$ :
les deux congruences valent aussi modulo $p$, d'où le résultat pour
$k\in\{1,2\}$.
\end{corrige}

\begin{pieges}
\begin{itemize}[leftmargin=1.4em, itemsep=2pt]
\item Ne pas voir que $Q^{\T}Q=I_n$ règle instantanément le premier
point : certains candidats se lancent dans des calculs de
coefficients inutiles.
\item Oublier d'invoquer l'hypothèse $Q^{\T}AQ\in\Mn(\Z)$ : sans elle,
$\ell^{2}Q^{\T}AQ$ n'a aucune raison d'avoir des coefficients dans
$p^{2}\Z$ (les coefficients de $Q^{\T}AQ$ pourraient avoir des
dénominateurs).
\item Traiter laborieusement $k=1$ puis $k=2$ : le cas $k=2$ implique
le cas $k=1$, il suffit de le dire.
\end{itemize}
\end{pieges}

% ------------------------------------------------ Q6
\begin{enonce}{Question 6}
Soit $U\in\Mn(\Z)$.
\begin{enumerate}
\item Si $U\in\GLn(\Z)$, démontrer que $\det(U)=\pm 1$.
\item Si $\det(U)=\pm 1$, démontrer que $U\in\GLn(\Z)$.
\item Montrer que $(\GLn(\Z),\times)$ est un groupe.
\end{enumerate}
\end{enonce}

\begin{corrige}
\textbf{(a)} Par définition, $U$ est inversible dans $\Mn(\R)$ et
$U^{-1}\in\Mn(\Z)$. Le déterminant d'une matrice à coefficients
entiers est un entier (formule de Leibniz) : $\det(U)\in\Z$ et
$\det(U^{-1})\in\Z$. Or
\[
\det(U)\,\det(U^{-1})=\det(UU^{-1})=\det(I_n)=1 .
\]
Ainsi $\det(U)$ est un élément inversible de l'anneau $\Z$, donc
$\det(U)\in\{-1,1\}$.

\textbf{(b)} Supposons $\det(U)=\varepsilon\in\{-1,1\}$. En
particulier $\det(U)\neq 0$, donc $U$ est inversible dans $\Mn(\R)$
et la formule de la comatrice donne
\[
U^{-1}=\frac{1}{\det(U)}\,{}\mathrm{com}(U)^{\T}
=\varepsilon\,\mathrm{com}(U)^{\T}.
\]
Chaque cofacteur de $U$ est, au signe près, le déterminant d'une
matrice extraite de $U$, donc à coefficients entiers : c'est un
entier. Ainsi $\mathrm{com}(U)^{\T}\in\Mn(\Z)$, puis
$U^{-1}=\varepsilon\,\mathrm{com}(U)^{\T}\in\Mn(\Z)$, et
$U\in\GLn(\Z)$.

\textbf{(c)} Montrons que $\GLn(\Z)$ est un \emph{sous-groupe} de
$(\GLn(\R),\times)$, ce qui suffira (l'associativité, l'élément
neutre et les inverses sont alors hérités).
\begin{itemize}[leftmargin=1.6em, itemsep=1pt]
\item $\GLn(\Z)\subset\GLn(\R)$ et $I_n\in\GLn(\Z)$ car
$I_n\in\Mn(\Z)$, $I_n^{-1}=I_n\in\Mn(\Z)$ : la partie est non vide.
\item \emph{Stabilité par produit :} si $U,V\in\GLn(\Z)$, alors
$UV\in\Mn(\Z)$ (les coefficients d'un produit de matrices entières
sont des sommes de produits d'entiers) et
$(UV)^{-1}=V^{-1}U^{-1}\in\Mn(\Z)$ pour la même raison. Donc
$UV\in\GLn(\Z)$.
\item \emph{Stabilité par inverse :} si $U\in\GLn(\Z)$, alors
$U^{-1}\in\Mn(\Z)$ par définition, et $(U^{-1})^{-1}=U\in\Mn(\Z)$ ;
donc $U^{-1}\in\GLn(\Z)$.
\end{itemize}
Ainsi $(\GLn(\Z),\times)$ est un groupe.
\end{corrige}

\begin{pieges}
\begin{itemize}[leftmargin=1.4em, itemsep=2pt]
\item En (a), écrire \og $\det(U^{-1})=1/\det(U)$ donc c'est fini \fg{}
sans utiliser que $\det(U^{-1})$ est \emph{entier} : le point clé est
que $\det(U)$ divise $1$ \emph{dans $\Z$}.
\item En (b), oublier de justifier que la comatrice est à coefficients
entiers, ou oublier de dire que $\det(U)\neq 0$ assure d'abord
l'inversibilité dans $\Mn(\R)$.
\item En (c), redémontrer l'associativité à la main au lieu de
travailler par sous-groupe ; ou, à l'inverse, oublier de vérifier que
le \emph{produit} reste à coefficients entiers (la structure demande
une loi \emph{interne}).
\end{itemize}
\end{pieges}

% ------------------------------------------------ Q7
\begin{enonce}{Question 7}
\begin{enumerate}
\item Soit $M\in\On(\Z)$. Démontrer que sur chaque ligne et chaque
colonne de $M$ il y a exactement un coefficient non nul, qui vaut
$\pm 1$.
\item En déduire le cardinal de $\On(\Z)$.
\end{enumerate}
\end{enonce}

\begin{corrige}
\textbf{(a)} Notons $M=[m_{i,j}]$. La relation $M^{\T}M=I_n$ signifie
que pour tous $j,j'$ :
\[
\sum_{i=1}^{n}m_{i,j}\,m_{i,j'}=\delta_{j,j'} .
\]
\emph{Colonnes.} Pour $j'=j$ : $\sum_{i}m_{i,j}^{2}=1$. C'est une
somme d'entiers positifs ou nuls valant $1$ : exactement un des
$m_{i,j}^{2}$ vaut $1$ (les autres étant nuls). Il existe donc un
unique indice, noté $\tau(j)$, tel que $m_{\tau(j),j}=\pm 1$, et
$m_{i,j}=0$ pour $i\neq\tau(j)$.

\emph{Lignes.} Pour $j\neq j'$, l'orthogonalité donne
$\sum_i m_{i,j}m_{i,j'}=0$. Si l'on avait $\tau(j)=\tau(j')$, cette
somme vaudrait $m_{\tau(j),j}\,m_{\tau(j),j'}=\pm 1\neq 0$ :
contradiction. Donc $\tau:\ii{1}{n}\to\ii{1}{n}$ est injective, et,
entre ensembles finis de même cardinal, bijective : c'est une
permutation. Chaque ligne $i$ contient alors exactement un coefficient
non nul, celui de la colonne $\tau^{-1}(i)$, qui vaut $\pm 1$.

\textbf{(b)} D'après (a), toute matrice $M\in\On(\Z)$ s'écrit de
manière unique
\[
M=[\,\varepsilon_j\,\delta_{i,\tau(j)}\,]_{1\le i,j\le n},
\qquad \tau\in\Perm,\quad
(\varepsilon_1,\dots,\varepsilon_n)\in\{-1,1\}^{n},
\]
où $\tau$ est la permutation ci-dessus et $\varepsilon_j=m_{\tau(j),j}$.
Réciproquement, toute matrice de cette forme est dans $\On(\Z)$ : ses
coefficients sont entiers, et ses colonnes sont, au signe près, des
vecteurs distincts de la base canonique, donc forment une famille
orthonormale : $M^{\T}M=I_n$. L'application
\[
\Perm\times\{-1,1\}^{n}\longrightarrow\On(\Z),\qquad
(\tau,\varepsilon)\longmapsto[\,\varepsilon_j\,\delta_{i,\tau(j)}\,]
\]
est donc une bijection, et
\[
\boxed{\ \#\On(\Z)=2^{\,n}\,n!\ }
\]
\end{corrige}

\begin{pieges}
\begin{itemize}[leftmargin=1.4em, itemsep=2pt]
\item Ne traiter que les colonnes et \og par symétrie les lignes \fg{}
en utilisant $MM^{\T}=I_n$ sans le justifier : $M^{\T}M=I_n$ donne
$M$ inversible d'inverse $M^{\T}$, \emph{puis} $MM^{\T}=I_n$ ; il
faut l'écrire (ou utiliser l'argument d'injectivité de $\tau$
ci-dessus, qui l'évite).
\item L'argument arithmétique clé est \og une somme de carrés
d'entiers vaut $1$ \fg{} : il doit apparaître explicitement. Sur
$\R$, la conclusion serait évidemment fausse !
\item En (b), donner $2^{n}n!$ sans vérifier la \emph{réciproque}
(toutes ces matrices sont orthogonales) ni l'\emph{injectivité} du
paramétrage.
\end{itemize}
\end{pieges}

% ------------------------------------------------ Q8
\begin{enonce}{Question 8}
Soit $\sigma\in\Perm$. On pose
$Q_\sigma=[\delta_{i,\sigma(j)}]_{1\le i,j\le n}$. Démontrer que
$Q_\sigma^{\T}=Q_{\sigma^{-1}}$, puis que $Q_\sigma\in\On(\Z)$.
\end{enonce}

\begin{corrige}
\emph{Transposée.} Pour tous $i,j\in\ii{1}{n}$ :
\[
\big(Q_\sigma^{\T}\big)_{i,j}
=\big(Q_\sigma\big)_{j,i}
=\delta_{j,\sigma(i)} .
\]
Or $j=\sigma(i)\iff i=\sigma^{-1}(j)$, de sorte que
$\delta_{j,\sigma(i)}=\delta_{i,\sigma^{-1}(j)}
=\big(Q_{\sigma^{-1}}\big)_{i,j}$. D'où
$Q_\sigma^{\T}=Q_{\sigma^{-1}}$.

\emph{Morphisme.} Montrons que pour $\sigma,\tau\in\Perm$ on a
$Q_\sigma Q_\tau=Q_{\sigma\circ\tau}$. En effet :
\[
\big(Q_\sigma Q_\tau\big)_{i,j}
=\sum_{k=1}^{n}\delta_{i,\sigma(k)}\,\delta_{k,\tau(j)}
=\delta_{i,\sigma(\tau(j))}
=\big(Q_{\sigma\circ\tau}\big)_{i,j},
\]
le seul terme éventuellement non nul de la somme étant celui d'indice
$k=\tau(j)$.

\emph{Conclusion.} $Q_\sigma$ est à coefficients dans $\{0,1\}
\subset\Z$, et
\[
Q_\sigma^{\T}\,Q_\sigma
=Q_{\sigma^{-1}}\,Q_\sigma
=Q_{\sigma^{-1}\circ\,\sigma}
=Q_{\mathrm{id}}
=I_n ,
\]
car $Q_{\mathrm{id}}=[\delta_{i,j}]=I_n$. Donc
$Q_\sigma\in\On(\Z)$.
\end{corrige}

\begin{pieges}
\begin{itemize}[leftmargin=1.4em, itemsep=2pt]
\item Erreurs d'indices : avec la convention de l'énoncé, la colonne
$j$ de $Q_\sigma$ porte son $1$ en \emph{ligne} $\sigma(j)$ ; toute
confusion ligne/colonne inverse $\sigma$ et $\sigma^{-1}$.
\item Dans le calcul du produit, ne pas identifier l'unique indice $k$
qui contribue ($k=\tau(j)$) : les manipulations de Kronecker doivent
être justifiées.
\item Conclure \og orthogonale donc dans $\On(\Z)$ \fg{} en oubliant de
mentionner que les coefficients sont entiers.
\end{itemize}
\end{pieges}

% ------------------------------------------------ Q9
\begin{enonce}{Question 9}
Pour $i\neq j$ dans $\ii{1}{n}$, $P_{i,j}$ désigne la matrice égale à
$I_n$ sauf : $1$ en positions $(i,j)$ et $(j,i)$, $0$ en positions
$(i,i)$ et $(j,j)$. Soit $A=[a_{k,l}]\in\Mn(\R)$.
\begin{enumerate}
\item Démontrer qu'il existe $\sigma\in\Perm$ telle que
$P_{i,j}=Q_\sigma$.
\item Démontrer soigneusement que la matrice obtenue à partir de $A$
en échangeant les lignes $i$ et $j$ est $P_{i,j}A$.
\end{enumerate}
\end{enonce}

\begin{corrige}
\textbf{(a)} Soit $\tau=(i\ j)\in\Perm$ la transposition échangeant
$i$ et $j$. Vérifions coefficient par coefficient que
$P_{i,j}=Q_\tau$, c'est-à-dire $(P_{i,j})_{k,l}=\delta_{k,\tau(l)}$
pour tous $k,l$ :
\begin{itemize}[leftmargin=1.6em, itemsep=1pt]
\item si $l\notin\{i,j\}$ : $\tau(l)=l$ et
$\delta_{k,\tau(l)}=\delta_{k,l}$, soit un $1$ en position $(l,l)$ et
des $0$ ailleurs, conformément à la définition de $P_{i,j}$ ;
\item si $l=i$ : $\tau(i)=j$ et $\delta_{k,\tau(i)}=\delta_{k,j}$ :
la colonne $i$ porte un unique $1$ en ligne $j$, conformément au
$1$ en position $(j,i)$ et aux $0$ en $(i,i)$ et ailleurs ;
\item si $l=j$ : symétriquement, unique $1$ en position $(i,j)$.
\end{itemize}
Donc $P_{i,j}=Q_\tau$ avec $\sigma=\tau=(i\ j)$.

\textbf{(b)} Calculons le coefficient générique de $P_{i,j}A=Q_\tau A$.
Pour tous $k,l\in\ii{1}{n}$ :
\[
\big(P_{i,j}A\big)_{k,l}
=\sum_{m=1}^{n}\delta_{k,\tau(m)}\,a_{m,l} .
\]
L'unique valeur de $m$ pour laquelle $\tau(m)=k$ est
$m=\tau^{-1}(k)=\tau(k)$ (une transposition est involutive), donc
\[
\big(P_{i,j}A\big)_{k,l}=a_{\tau(k),\,l}\qquad
\text{pour tous }k,l .
\]
Autrement dit : la ligne $k$ de $P_{i,j}A$ est la ligne $\tau(k)$ de
$A$. Comme $\tau(i)=j$, $\tau(j)=i$ et $\tau(k)=k$ pour
$k\notin\{i,j\}$, la matrice $P_{i,j}A$ est exactement la matrice
obtenue à partir de $A$ en échangeant les lignes $i$ et $j$, les
autres lignes étant inchangées.
\end{corrige}

\begin{pieges}
\begin{itemize}[leftmargin=1.4em, itemsep=2pt]
\item L'adverbe \og soigneusement \fg{} est un signal du jury : un
schéma, un exemple $3\times 3$ ou un \og on voit que \fg{} ne
rapportent rien. Le calcul du coefficient générique avec le symbole
de Kronecker est attendu.
\item Oublier l'involutivité ($\tau^{-1}=\tau$) et écrire
$a_{\tau^{-1}(k),l}$ sans conclure.
\item Confondre le sens : \emph{multiplier à gauche agit sur les
lignes}, à droite sur les colonnes (l'énoncé admet d'ailleurs le cas
des colonnes).
\end{itemize}
\end{pieges}

% ------------------------------------------------ Q10
\begin{enonce}{Question 10}
Soient $q\in\Z$ et $i\neq j$ dans $\ii{1}{n}$. On pose
$T_{i,j}(q)=I_n-q\,E_{i,j}$.
\begin{enumerate}
\item Démontrer que $T_{i,j}(q)\in\GLn(\Z)$.
\item Démontrer soigneusement que la matrice obtenue à partir de $A$
par l'opération $C_j\leftarrow C_j-qC_i$ est $A\,T_{i,j}(q)$.
\item Donner, sans démonstration, $S_{i,j}(q)\in\GLn(\Z)$ telle que
$S_{i,j}(q)A$ réalise $L_j\leftarrow L_j-qL_i$.
\end{enumerate}
\end{enonce}

\begin{corrige}
\textbf{(a)} Rappelons la règle de multiplication des matrices
élémentaires : $E_{i,j}E_{k,l}=\delta_{j,k}\,E_{i,l}$. Comme
$i\neq j$, on a $E_{i,j}^{2}=\delta_{j,i}E_{i,j}=0$, d'où
\[
\big(I_n-qE_{i,j}\big)\big(I_n+qE_{i,j}\big)
=I_n-q^{2}E_{i,j}^{2}=I_n ,
\]
et de même dans l'autre sens. Ainsi $T_{i,j}(q)$ est inversible,
d'inverse $T_{i,j}(-q)=I_n+qE_{i,j}\in\Mn(\Z)$. Les deux matrices
étant à coefficients entiers, $T_{i,j}(q)\in\GLn(\Z)$.

\emph{Variante :} $T_{i,j}(q)$ est triangulaire à diagonale de $1$
(quitte à regarder le cas $i<j$ ou $i>j$), donc $\det=1$, et l'on
conclut par la question 6(b).

\textbf{(b)} On a $A\,T_{i,j}(q)=A-q\,A\,E_{i,j}$. Calculons le
coefficient générique de $AE_{i,j}$ : pour tous $k,l$,
\[
\big(AE_{i,j}\big)_{k,l}
=\sum_{m=1}^{n}a_{k,m}\,\big(E_{i,j}\big)_{m,l}
=\sum_{m=1}^{n}a_{k,m}\,\delta_{m,i}\,\delta_{l,j}
=a_{k,i}\,\delta_{l,j} .
\]
Ainsi $AE_{i,j}$ a toutes ses colonnes nulles sauf la $j$-ième, qui
est la colonne $C_i$ de $A$. Par conséquent, pour tous $k,l$ :
\[
\big(A\,T_{i,j}(q)\big)_{k,l}
=a_{k,l}-q\,a_{k,i}\,\delta_{l,j}
=\begin{cases}
a_{k,l} & \text{si } l\neq j,\\[2pt]
a_{k,j}-q\,a_{k,i} & \text{si } l=j .
\end{cases}
\]
Les colonnes de $A\,T_{i,j}(q)$ sont donc celles de $A$, à
l'exception de la colonne $j$ qui est remplacée par $C_j-qC_i$ :
c'est exactement l'opération $C_j\leftarrow C_j-qC_i$.

\textbf{(c)} On peut prendre
\[
S_{i,j}(q)=I_n-q\,E_{j,i}\;=\;T_{j,i}(q)\in\GLn(\Z).
\]
\end{corrige}

\begin{pieges}
\begin{itemize}[leftmargin=1.4em, itemsep=2pt]
\item Affirmer \og triangulaire \fg{} sans préciser
supérieure/inférieure selon la position relative de $i$ et $j$ ; la
preuve par l'inverse explicite $T_{i,j}(-q)$ évite toute discussion.
\item Utiliser $E_{i,j}^{2}=0$ sans mentionner l'hypothèse $i\neq j$
(pour $i=j$, $E_{i,i}^2=E_{i,i}\neq 0$).
\item En (c), se tromper d'indices : l'opération sur la \emph{ligne}
$j$ demande $E_{j,i}$ (le $-q$ se trouve en position $(j,i)$), pas
$E_{i,j}$. Une vérification rapide sur le coefficient générique,
même non demandée, sécurise la réponse.
\end{itemize}
\end{pieges}

% =====================================================================
\section{Partie II. Résultant et applications}
% =====================================================================

\noindent\emph{Rappel du cadre.} Soient $m,n\in\N^{*}$ avec
$m\ge 2$, et $R=\sum_{k=0}^{m}r_kX^{k}$, $S=\sum_{k=0}^{n}s_kX^{k}$
dans $\K[X]$, avec $\dg{R}=m$ \emph{exactement} et $\dg{S}\le n$. On
pose
\[
\varphi_{R,S}:\left\{
\begin{array}{ccl}
\K_{n-1}[X]\times\K_{m-1}[X] & \longrightarrow & \K_{n+m-1}[X]\\
(G,H) & \longmapsto & GR+HS
\end{array}\right.
\]
puis $D=R\wedge S$, $d=\dg{D}$, et
$\psi_R=\varphi_{R,R'}:\K_{m-2}[X]\times\K_{m-1}[X]\to\K_{2m-2}[X]$,
$\Delta(R)=\det(\psi_R)$.

% ------------------------------------------------ Q11
\begin{enonce}{Question 11}
Démontrer le lemme de Gauss : si $G,H,L\in\K[X]$ vérifient
$G\wedge H=1$ et $G\mid HL$, alors $G\mid L$.
\end{enonce}

\begin{corrige}
L'anneau $\K[X]$ est euclidien (pour le degré), donc principal ; le
théorème de Bézout s'y applique : puisque $G\wedge H=1$, il existe
$U,V\in\K[X]$ tels que
\[
UG+VH=1 .
\]
Multiplions cette relation par $L$ :
\[
L=UGL+VHL .
\]
Le polynôme $G$ divise $UGL$ (évident) ; il divise $HL$ par
hypothèse, donc aussi $VHL$. Par somme, $G$ divise
$UGL+VHL=L$, ce qu'il fallait démontrer. Précisément, si $HL=GM$
avec $M\in\K[X]$, alors $L=G\,(UL+VM)$.
\end{corrige}

\begin{pieges}
\begin{itemize}[leftmargin=1.4em, itemsep=2pt]
\item Invoquer Bézout sans dire pourquoi il est valable ($\K[X]$
euclidien donc principal) : c'est précisément le contenu de la
question, qui est un exercice de cours à rédiger sans trou.
\item Raisonner par décomposition en facteurs irréductibles : c'est
correct dans l'anneau factoriel $\K[X]$, mais la rédaction en est
plus délicate (unicité de la décomposition à invoquer proprement) ;
la voie Bézout est la voie attendue.
\item Attention : ce lemme, vrai ici grâce à Bézout, sert ensuite en
question 15 ; il faut le citer à ce moment-là.
\end{itemize}
\end{pieges}

% ------------------------------------------------ Q12
\begin{enonce}{Question 12}
Démontrer que $\varphi_{R,S}$ est une application linéaire de
$\K_{n-1}[X]\times\K_{m-1}[X]$ dans $\K_{n+m-1}[X]$.
\end{enonce}

\begin{corrige}
\emph{L'application est bien à valeurs dans $\K_{n+m-1}[X]$.} Soit
$(G,H)\in\K_{n-1}[X]\times\K_{m-1}[X]$. Alors
\[
\dg{GR}\le\dg{G}+\dg{R}\le (n-1)+m=n+m-1,
\qquad
\dg{HS}\le (m-1)+n=n+m-1,
\]
avec les conventions usuelles sur le degré du polynôme nul. Donc
$GR+HS\in\K_{n+m-1}[X]$.

\emph{Linéarité.} Soient $(G,H),(G',H')$ dans l'espace de départ et
$\lambda\in\K$. Alors, par distributivité dans $\K[X]$ :
\begin{align*}
\varphi_{R,S}\big(\lambda(G,H)+(G',H')\big)
&=\varphi_{R,S}\big(\lambda G+G',\,\lambda H+H'\big)\\
&=(\lambda G+G')R+(\lambda H+H')S\\
&=\lambda\,(GR+HS)+(G'R+H'S)\\
&=\lambda\,\varphi_{R,S}(G,H)+\varphi_{R,S}(G',H').
\end{align*}
Donc $\varphi_{R,S}$ est linéaire.
\end{corrige}

\begin{pieges}
\begin{itemize}[leftmargin=1.4em, itemsep=2pt]
\item Oublier la vérification des degrés : montrer qu'une application
est \og linéaire de $E$ dans $F$ \fg{} inclut le fait qu'elle est
\emph{bien définie à valeurs dans $F$}.
\item Confondre la linéarité en $(G,H)$ (demandée) avec une
prétendue linéarité en $(R,S)$ (qui sont fixés).
\item La structure d'espace vectoriel produit doit être utilisée
correctement : $\lambda(G,H)=(\lambda G,\lambda H)$.
\end{itemize}
\end{pieges}

% ------------------------------------------------ Q13
\begin{enonce}{Question 13}
Soit
$\mathcal{U}=\big((1,0),(X,0),\dots,(X^{n-1},0),(0,1),(0,X),\dots,
(0,X^{m-1})\big)$. Démontrer que $\mathcal{U}$ est une base de
$\K_{n-1}[X]\times\K_{m-1}[X]$.
\end{enonce}

\begin{corrige}
L'espace produit de deux espaces de dimension finie est de dimension
la somme des dimensions :
\[
\dim\big(\K_{n-1}[X]\times\K_{m-1}[X]\big)
=\dim\K_{n-1}[X]+\dim\K_{m-1}[X]=n+m .
\]
La famille $\mathcal{U}$ compte exactement $n+m$ vecteurs ; il suffit
donc de montrer qu'elle est \emph{libre}.

Soient $a_0,\dots,a_{n-1},b_0,\dots,b_{m-1}\in\K$ tels que
\[
\sum_{k=0}^{n-1}a_k\,(X^{k},0)+\sum_{k=0}^{m-1}b_k\,(0,X^{k})=(0,0).
\]
Par définition des lois du produit, cette égalité s'écrit
\[
\Big(\sum_{k=0}^{n-1}a_kX^{k},\ \sum_{k=0}^{m-1}b_kX^{k}\Big)=(0,0),
\]
soit $\sum_k a_kX^k=0$ et $\sum_k b_kX^k=0$. Comme
$(1,X,\dots,X^{n-1})$ et $(1,X,\dots,X^{m-1})$ sont les bases
canoniques (donc des familles libres) de $\K_{n-1}[X]$ et
$\K_{m-1}[X]$, tous les $a_k$ et tous les $b_k$ sont nuls.

La famille $\mathcal{U}$ est libre de cardinal $n+m=\dim$, c'est une
base.
\end{corrige}

\begin{pieges}
\begin{itemize}[leftmargin=1.4em, itemsep=2pt]
\item Décréter \og c'est la base canonique du produit \fg{} : c'est
justement ce qu'on demande de prouver (et le Vrai/Faux 2 vient de
montrer qu'on peut se tromper sur les bases d'un produit !).
\item Oublier de citer la formule
$\dim(E\times F)=\dim E+\dim F$, ou compter $nm$ vecteurs.
\item Montrer le caractère générateur \emph{et} la liberté alors que
l'argument \og libre + bon cardinal \fg{} économise la moitié du
travail.
\end{itemize}
\end{pieges}

% ------------------------------------------------ Q14
\begin{enonce}{Question 14}
Expliciter les coefficients de
$\Mat_{\mathcal{U},\mathcal{V}}(\varphi_{R,S})$, où
$\mathcal{V}=(1,X,\dots,X^{m+n-1})$ est la base canonique de
$\K_{n+m-1}[X]$.
\end{enonce}

\begin{corrige}
Adoptons les conventions $r_k=0$ pour $k\notin\ii{0}{m}$ et
$s_k=0$ pour $k\notin\ii{0}{n}$.

Les colonnes de la matrice sont les coordonnées dans $\mathcal{V}$
des images des vecteurs de $\mathcal{U}$. Pour
$j\in\ii{1}{n}$, le $j$-ième vecteur de $\mathcal{U}$ est
$(X^{j-1},0)$ et
\[
\varphi_{R,S}(X^{j-1},0)=X^{j-1}R
=\sum_{k=0}^{m}r_k\,X^{k+j-1}
=\sum_{i=1}^{n+m}r_{\,i-j}\;X^{\,i-1},
\]
en posant $i=k+j$ (le coefficient de $X^{i-1}$ est $r_{i-j}$). De
même, pour $j\in\ii{1}{m}$, le $(n+j)$-ième vecteur de $\mathcal{U}$
est $(0,X^{j-1})$ et
\[
\varphi_{R,S}(0,X^{j-1})=X^{j-1}S
=\sum_{i=1}^{n+m}s_{\,i-j}\;X^{\,i-1}.
\]
Notant $\Sigma=\Mat_{\mathcal{U},\mathcal{V}}(\varphi_{R,S})
\in\mathcal{M}_{n+m}(\K)$, on a donc, pour tout
$i\in\ii{1}{n+m}$ :
\[
\boxed{\;\Sigma_{i,j}=r_{\,i-j}\ \ (1\le j\le n),
\qquad
\Sigma_{i,\,n+j}=s_{\,i-j}\ \ (1\le j\le m).\;}
\]
Chaque colonne du premier bloc contient la liste
$r_0,r_1,\dots,r_m$ décalée d'un cran vers le bas à chaque colonne ;
idem pour $s_0,\dots,s_n$ dans le second bloc :
\[
\Sigma=
\begin{pmatrix}
r_0    &        &        &(0)     & s_0    &        &  &(0)  \\
r_1    & r_0    &        &        & s_1    & s_0    &  &     \\
\vdots & r_1    & \ddots &        & \vdots & s_1    & \ddots & \\
r_m    & \vdots & \ddots & r_0    & s_n    & \vdots & \ddots & s_0\\
       & r_m    &        & r_1    &        & s_n    &        & s_1\\
       &        & \ddots & \vdots &        &        & \ddots & \vdots\\
(0)    &        &        & r_m    & (0)    &        &        & s_n
\end{pmatrix}
\begin{matrix}
\left.\vphantom{\begin{matrix}1\\1\\1\\1\\1\\1\\1\end{matrix}}
\right\}\ n+m\ \text{lignes}
\end{matrix}
\]
(le premier bloc a $n$ colonnes, le second $m$). C'est la
\emph{matrice de Sylvester} de $(R,S)$ ; son déterminant, dans le cas
$S=R'$, est $\Delta(R)$.
\end{corrige}

\begin{pieges}
\begin{itemize}[leftmargin=1.4em, itemsep=2pt]
\item L'erreur classique : rendre la matrice \emph{transposée}. Ici les
\emph{colonnes} sont les images des vecteurs de $\mathcal{U}$ ; la
matrice de Sylvester est souvent présentée en lignes dans les
ouvrages, d'où la confusion.
\item Oublier les conventions $r_k=0$ hors de $\ii{0}{m}$ : sans
elles, la formule $\Sigma_{i,j}=r_{i-j}$ n'a pas de sens pour tous
les indices.
\item Décalages d'indices entre \og coefficient de $X^{i-1}$ \fg{} et
\og $i$-ième ligne \fg{} : poser explicitement le changement d'indice
$i=k+j$ sécurise le calcul.
\end{itemize}
\end{pieges}

% ------------------------------------------------ Q15
\begin{enonce}{Question 15}
Soient $R_1,S_1\in\K[X]$ tels que $R=DR_1$ et $S=DS_1$. Démontrer que
\[
\Ker(\varphi_{R,S})=\big\{(S_1W,\,-R_1W),\ W\in\K_{d-1}[X]\big\},
\]
et en déduire $\dim\Ker(\varphi_{R,S})$.
\end{enonce}

\begin{corrige}
Remarquons d'abord que $R\neq 0$ (car $\dg{R}=m$), donc
$D=R\wedge S\neq 0$ et $R_1\neq 0$, avec
\[
\dg{R_1}=m-d\qquad\text{et}\qquad \dg{S_1}\le n-d .
\]
De plus $R_1\wedge S_1=1$ : si un polynôme irréductible $P$ divisait
$R_1$ et $S_1$, alors $DP$ diviserait $R$ et $S$, donc $DP\mid D$,
ce qui est impossible pour des raisons de degré. (Si $S=0$, alors
$S_1=0$ et la convention $R_1\wedge 0=R_1$, avec $D=R$ unitaire
associé à $R$, donne $R_1$ constant non nul : le raisonnement
ci-dessous couvre ce cas.)

\emph{Inclusion $\supseteq$.} Soit $W\in\K_{d-1}[X]$. Vérifions
d'abord que le couple appartient bien à l'espace de départ :
\[
\dg{S_1W}\le (n-d)+(d-1)=n-1,
\qquad
\dg{R_1W}\le (m-d)+(d-1)=m-1 ,
\]
donc $(S_1W,-R_1W)\in\K_{n-1}[X]\times\K_{m-1}[X]$. Ensuite :
\[
\varphi_{R,S}(S_1W,-R_1W)
=S_1W\,R-R_1W\,S
=S_1W\,DR_1-R_1W\,DS_1
=0 .
\]

\emph{Inclusion $\subseteq$.} Soit
$(G,H)\in\Ker(\varphi_{R,S})$ : $GR+HS=0$, soit
$G\,DR_1=-H\,DS_1$. L'anneau $\K[X]$ étant intègre et $D\neq 0$, on
simplifie par $D$ :
\[
G\,R_1=-H\,S_1 . \tag{$*$}
\]
\begin{itemize}[leftmargin=1.6em, itemsep=1pt]
\item \emph{Si $S_1=0$} (cas $S=0$) : $(*)$ donne $GR_1=0$ avec
$R_1\neq 0$, donc $G=0$ ; ici $d=m$, $R_1=r_m$ (constante non
nulle), et tout $H\in\K_{m-1}[X]$ convient : en posant
$W=-H/r_m\in\K_{d-1}[X]$, on a bien $(G,H)=(S_1W,-R_1W)$.
\item \emph{Si $S_1\neq 0$} : de $(*)$, $S_1$ divise $G R_1$ ; comme
$S_1\wedge R_1=1$, le lemme de Gauss (question 11) donne
$S_1\mid G$ : il existe $W\in\K[X]$ tel que $G=S_1W$. En reportant
dans $(*)$ : $S_1WR_1=-HS_1$, et en simplifiant par $S_1\neq 0$
(intégrité) : $H=-R_1W$.

\emph{Contrôle du degré de $W$.} Comme $\dg{H}\le m-1$ et
$\dg{R_1}=m-d$ \emph{exactement}, on a, si $W\neq 0$,
\[
\dg{W}=\dg{H}-\dg{R_1}\le (m-1)-(m-d)=d-1,
\]
donc $W\in\K_{d-1}[X]$.
\end{itemize}
Dans tous les cas $(G,H)=(S_1W,-R_1W)$ avec $W\in\K_{d-1}[X]$, ce
qui achève la double inclusion.

\emph{Dimension.} L'application
\[
\Theta:\K_{d-1}[X]\longrightarrow\Ker(\varphi_{R,S}),
\qquad W\longmapsto(S_1W,-R_1W)
\]
est linéaire, surjective (ce qui précède) et injective : si
$(S_1W,-R_1W)=(0,0)$, alors $R_1W=0$ avec $R_1\neq 0$, donc $W=0$.
C'est un isomorphisme, d'où
\[
\boxed{\ \dim\Ker(\varphi_{R,S})=\dim\K_{d-1}[X]=d
=\dg{R\wedge S}.\ }
\]
(Si $d=0$, $\K_{d-1}[X]=\{0\}$ et $\varphi_{R,S}$ est injective.)
\end{corrige}

\begin{pieges}
\begin{itemize}[leftmargin=1.4em, itemsep=2pt]
\item \textbf{Oublier de justifier $R_1\wedge S_1=1$} avant d'invoquer
Gauss : c'est le point central, et il se démontre en une ligne par
maximalité de $D$.
\item \textbf{Négliger les contraintes de degré} : dans l'inclusion
$\supseteq$, vérifier que $(S_1W,-R_1W)$ est dans l'espace de
\emph{départ} ; dans $\subseteq$, borner $\dg{W}$ en utilisant que
$\dg{R_1}=m-d$ \emph{exactement} (côté $S_1$, le degré peut être
strictement inférieur à $n-d$, la majoration par $H$ est la bonne).
\item Simplifier par $D$ ou par $S_1$ sans mentionner l'intégrité de
$\K[X]$, ou sans avoir écarté le cas $S_1=0$.
\item Conclure $\dim=d-1$ : la dimension de $\K_{d-1}[X]$ est $d$
(degré $\le d-1$, donc $d$ coefficients).
\end{itemize}
\end{pieges}

% ------------------------------------------------ Q16
\begin{enonce}{Question 16}
Déterminer $\Ima(\varphi_{R,S})$.
\end{enonce}

\begin{corrige}
Montrons que
\[
\boxed{\ \Ima(\varphi_{R,S})
=D\cdot\K_{n+m-1-d}[X]
=\big\{DP,\ P\in\K[X],\ \dg{P}\le n+m-1-d\big\},\ }
\]
c'est-à-dire l'ensemble des multiples de $D$ appartenant à
$\K_{n+m-1}[X]$.

\emph{Inclusion directe.} Pour $(G,H)$ dans l'espace de départ,
\[
GR+HS=D\,(GR_1+HS_1)
\]
est un multiple de $D$, de degré $\le n+m-1$ ; en écrivant
$GR+HS=DP$, l'intégrité donne $\dg{P}=\dg{GR+HS}-d\le n+m-1-d$
lorsque $GR+HS\neq 0$ (et $0=D\cdot 0$ convient). D'où
$\Ima(\varphi_{R,S})\subseteq D\cdot\K_{n+m-1-d}[X]$.

\emph{Égalité par les dimensions.} D'une part, le théorème du rang
appliqué à $\varphi_{R,S}$ et la question 15 donnent
\[
\dim\Ima(\varphi_{R,S})
=\dim\big(\K_{n-1}[X]\times\K_{m-1}[X]\big)-\dim\Ker(\varphi_{R,S})
=(n+m)-d .
\]
D'autre part, l'application $P\mapsto DP$ est linéaire injective
(intégrité, $D\neq 0$) de $\K_{n+m-1-d}[X]$ sur
$D\cdot\K_{n+m-1-d}[X]$, donc
\[
\dim\big(D\cdot\K_{n+m-1-d}[X]\big)=n+m-d .
\]
Une inclusion entre sous-espaces de même dimension finie est une
égalité : $\Ima(\varphi_{R,S})=D\cdot\K_{n+m-1-d}[X]$.

\emph{Cas particulier utile.} Si $R\wedge S=1$ (soit $d=0$),
$\varphi_{R,S}$ est un isomorphisme sur $\K_{n+m-1}[X]$ : c'est
l'identité de Bézout \og avec contrôle des degrés \fg.
\end{corrige}

\begin{pieges}
\begin{itemize}[leftmargin=1.4em, itemsep=2pt]
\item Tenter la double inclusion \og à la main \fg{} : l'inclusion
réciproque directe (tout multiple de $D$ de bon degré est atteint)
est faisable mais délicate ; le théorème du rang, combiné à la
question 15, donne l'égalité en trois lignes. Savoir \emph{choisir}
l'argument fait partie de l'épreuve.
\item Erreur de dimension : $D\cdot\K_{n+m-1-d}[X]$ est de dimension
$n+m-d$ (et non $n+m-d-1$) ; l'isomorphisme $P\mapsto DP$ doit être
mentionné.
\item Répondre \og $\Ima=\K_{n+m-1}[X]$ \fg{} : ce n'est vrai que si
$d=0$.
\end{itemize}
\end{pieges}

% ------------------------------------------------ Q17
\begin{enonce}{Question 17}
Démontrer que si $R$ a un facteur multiple dans sa décomposition en
facteurs irréductibles, alors $\Delta(R)=0$.
\end{enonce}

\begin{corrige}
Par hypothèse, il existe $P\in\K[X]$ irréductible et $T\in\K[X]$
tels que $R=P^{2}T$. Dérivons :
\[
R'=2PP'T+P^{2}T'=P\,\big(2P'T+PT'\big),
\]
de sorte que $P$ divise à la fois $R$ et $R'$. Ainsi $P$ divise
$D_\psi:=R\wedge R'$, et donc
\[
\dg{R\wedge R'}\ \ge\ \dg{P}\ \ge\ 1 .
\]
Or $\psi_R=\varphi_{R,R'}$ entre dans le cadre de la question 15
appliqué au couple $(R,S)=(R,R')$, avec ici $\dg{R'}\le m-1$ (le
paramètre \og $n$ \fg{} vaut $m-1$) : d'après 15,
\[
\dim\Ker(\psi_R)=\dg{R\wedge R'}\ \ge\ 1 .
\]
L'application linéaire $\psi_R$ relie deux espaces de même dimension
finie,
\[
\dim\big(\K_{m-2}[X]\times\K_{m-1}[X]\big)=(m-1)+m=2m-1
=\dim\K_{2m-2}[X],
\]
et son noyau est non trivial : elle n'est pas injective, donc pas
bijective, et sa matrice dans les bases $\mathcal{U},\mathcal{V}$
est non inversible. Par conséquent
\[
\Delta(R)=\det(\psi_R)=0 .
\]
\emph{Remarque (caractéristique quelconque).} L'argument n'utilise
jamais $2\neq 0$ ni $R'\neq 0$ : même si $R'=0$ (possible en
caractéristique $p$, cf. Vrai/Faux 4), on a $R\wedge R'=R\wedge 0$
associé à $R$, de degré $m\ge 1$, et la conclusion subsiste.
\end{corrige}

\begin{pieges}
\begin{itemize}[leftmargin=1.4em, itemsep=2pt]
\item Raisonner avec une \og racine commune de $R$ et $R'$ \fg{} : un
facteur irréductible n'a pas nécessairement de racine dans $\K$
(penser à $X^2+1$ sur $\R$). Il faut raisonner en termes de
\emph{divisibilité par $P$}.
\item Écrire $R'=2PP'T+\dots$ puis simplifier par $2$ : inutile ici,
mais surtout illicite en caractéristique $2$ ; la factorisation par
$P$ suffit.
\item Oublier de vérifier que départ et arrivée de $\psi_R$ ont même
dimension $2m-1$ : parler de $\det(\psi_R)$ n'a de sens que pour une
matrice \emph{carrée}.
\end{itemize}
\end{pieges}

% ------------------------------------------------ Q18
\begin{enonce}{Question 18}
On suppose $\dim\Ker(\psi_R)=1$.
\begin{enumerate}
\item Démontrer qu'il existe $\lambda\in\K$ et $T\in\K[X]$ tels que
$T(\lambda)\neq 0$ et $R=(X-\lambda)^{2}T$.
\item Démontrer qu'alors
$\Ker(\psi_R)=\vect\big((2T+(X-\lambda)T',\,-(X-\lambda)T)\big)$.
\end{enumerate}
\end{enonce}

\begin{corrige}
\textbf{(a)} D'après la question 15 appliquée à $(R,R')$,
\[
\dg{R\wedge R'}=\dim\Ker(\psi_R)=1 .
\]
Le polynôme $R\wedge R'$ est unitaire de degré $1$ : il existe
$\lambda\in\K$ tel que
\[
R\wedge R'=X-\lambda .
\]
En particulier $(X-\lambda)\mid R$ : notons $k\ge 1$ la multiplicité
de $\lambda$ dans $R$, c'est-à-dire $R=(X-\lambda)^{k}T$ avec
$T\in\K[X]$ et $T(\lambda)\neq 0$. Dérivons :
\[
R'=(X-\lambda)^{k-1}\big(\,\underbrace{k}_{\in\K}\,T
+(X-\lambda)T'\big).
\]
\begin{itemize}[leftmargin=1.6em, itemsep=1pt]
\item \emph{Si $k=1$} : $R'=T+(X-\lambda)T'$, d'où
$R'(\lambda)=T(\lambda)\neq 0$, ce qui contredit
$(X-\lambda)\mid R'$ (conséquence de $R\wedge R'=X-\lambda$).
\item \emph{Si $k\ge 3$} : $(X-\lambda)^{k-1}$ divise $R'$ avec
$k-1\ge 2$, et $(X-\lambda)^{k}$ divise $R$ ; donc
$(X-\lambda)^{2}$ divise $R\wedge R'$, ce qui contredit
$\dg{R\wedge R'}=1$.
\end{itemize}
Il reste $k=2$ :
$R=(X-\lambda)^{2}T$ avec $T(\lambda)\neq 0$.

\textbf{(b)} Appliquons la description de la question 15 au couple
$(R,S)=(R,R')$, dont le pgcd est $D=X-\lambda$, soit $d=1$. On
factorise :
\[
R=(X-\lambda)\cdot\underbrace{(X-\lambda)T}_{=:R_1},
\qquad
R'=2(X-\lambda)T+(X-\lambda)^{2}T'
=(X-\lambda)\cdot\underbrace{\big(2T+(X-\lambda)T'\big)}_{=:S_1}.
\]
La question 15 donne alors
\[
\Ker(\psi_R)
=\big\{(S_1W,\,-R_1W),\ W\in\K_{0}[X]\big\}
=\big\{\big(\alpha S_1,\,-\alpha R_1\big),\ \alpha\in\K\big\}
=\vect\big((S_1,-R_1)\big),
\]
c'est-à-dire
\[
\boxed{\ \Ker(\psi_R)
=\vect\Big(\big(2T+(X-\lambda)T',\ -(X-\lambda)T\big)\Big).\ }
\]
Le vecteur directeur est non nul (sa seconde composante
$-(X-\lambda)T$ est non nulle), en cohérence avec
$\dim\Ker(\psi_R)=1$.
\end{corrige}

\begin{pieges}
\begin{itemize}[leftmargin=1.4em, itemsep=2pt]
\item \textbf{Le piège de la caractéristique :} écrire
\og $R'=k(X-\lambda)^{k-1}T+\cdots$ et $k\neq 0$ donc\dots \fg{} est
illicite : $k$ désigne ici l'image de l'entier $k$ dans $\K$, qui
peut être nulle (par exemple $k=p$ en caractéristique $p$). La
discussion $k=1$ / $k\ge 3$ ci-dessus contourne complètement cet
écueil.
\item Utiliser la caractérisation \og $\lambda$ racine de multiplicité
$\ge 2$ $\iff$ $R(\lambda)=R'(\lambda)=0$ \fg{} sans précaution :
son sens réciproque usuel repose sur la formule de Taylor, invalide
en caractéristique positive.
\item En (b), oublier que $d=1$ force $W$ \emph{constant}
($\K_0[X]=\K$) : c'est ce qui fait apparaître une droite
vectorielle.
\item Erreurs de signe dans $(S_1W,-R_1W)$ : le couple est
$(\,\cdot\,,-\,\cdot\,)$, à reporter fidèlement.
\end{itemize}
\end{pieges}

% ------------------------------------------------ Q19
\begin{enonce}{Question 19}
Démontrer que s'il existe $T,W\in\K[X]$ avec $\dg{T}\ge 2$ tels que
$R=T^{2}W$, alors $\dim\Ker(\psi_R)\ge 2$.
\end{enonce}

\begin{corrige}
Dérivons $R=T^{2}W$ :
\[
R'=2TT'W+T^{2}W'=T\,\big(2T'W+TW'\big),
\]
de sorte que $T$ divise $R$ et $R'$, donc $T$ divise $R\wedge R'$.
Par croissance du degré le long de la divisibilité (les polynômes en
jeu sont non nuls : $R\neq 0$ donc $T\neq 0$),
\[
\dg{R\wedge R'}\ \ge\ \dg{T}\ \ge\ 2 .
\]
La question 15, appliquée comme précédemment au couple $(R,R')$,
donne alors
\[
\dim\Ker(\psi_R)=\dg{R\wedge R'}\ \ge\ 2 .
\]
\end{corrige}

\begin{pieges}
\begin{itemize}[leftmargin=1.4em, itemsep=2pt]
\item Chercher à exhiber \og à la main \fg{} deux éléments
linéairement indépendants du noyau : c'est possible (par exemple
$(T'W\,U,\;-\ldots)$ pour divers $U$) mais acrobatique ; la formule
$\dim\Ker(\psi_R)=\dg{R\wedge R'}$ de la question 15 rend la
question immédiate. Le sujet est conçu pour être \emph{réinvesti}.
\item Justifier $T\mid R\wedge R'$ en passant par les racines de $T$ :
même écueil qu'en question 17, $T$ peut n'avoir aucune racine dans
$\K$ ; la divisibilité se lit sur les factorisations.
\item Ne pas remarquer que $T\neq 0$ (sinon $R=0$, exclu car
$\dg{R}=m$) avant de comparer les degrés.
\end{itemize}
\end{pieges}

% ------------------------------------------------ Q20
\begin{enonce}{Question 20}
En déduire que si $\dim\Ker(\psi_R)=1$, alors il existe
$\lambda\in\K$, $\mu\in\K\setminus\{0\}$ et des polynômes
$T_1,\dots,T_r\in\K[X]$ deux à deux distincts et irréductibles tels
que
\[
R=\mu\,(X-\lambda)^{2}\prod_{i=1}^{r}T_i,
\qquad\text{avec }\forall i\in\ii{1}{r},\ T_i(\lambda)\neq 0 .
\]
\end{enonce}

\begin{corrige}
Supposons $\dim\Ker(\psi_R)=1$. D'après la question 18(a), il existe
$\lambda\in\K$ et $T\in\K[X]$ tels que
\[
R=(X-\lambda)^{2}\,T,\qquad T(\lambda)\neq 0 .
\]
L'anneau $\K[X]$ est factoriel : écrivons la décomposition de $T$ en
facteurs irréductibles,
\[
T=\mu\prod_{i=1}^{r}T_i^{\,\alpha_i},
\]
avec $\mu\in\K\setminus\{0\}$ (le coefficient dominant de $T$), les
$T_i$ irréductibles unitaires deux à deux distincts et
$\alpha_i\ge 1$ (le cas $r=0$, $T$ constant, étant autorisé avec un
produit vide).

\emph{Chaque $\alpha_i$ vaut $1$.} Supposons par l'absurde
$\alpha_{i_0}\ge 2$ pour un indice $i_0$. Alors
$T_{i_0}^{2}\mid T$, et l'on peut écrire
\[
R=(X-\lambda)^{2}\,T_{i_0}^{2}\,V
=\big(\,\underbrace{(X-\lambda)\,T_{i_0}}_{=:\widetilde T}\,\big)^{2}
\,V,\qquad V\in\K[X],
\]
avec $\dg{\widetilde T}=1+\dg{T_{i_0}}\ge 2$. La question 19
s'applique et donne $\dim\Ker(\psi_R)\ge 2$ : contradiction. Donc
$\alpha_i=1$ pour tout $i$.

\emph{Les $T_i$ ne s'annulent pas en $\lambda$.} Pour tout $i$,
$T_i$ divise $T$ ; si l'on avait $T_i(\lambda)=0$, alors
$T(\lambda)=0$, contredisant $T(\lambda)\neq 0$. Donc
$T_i(\lambda)\neq 0$ pour tout $i$ ; en particulier aucun $T_i$
n'est associé à $X-\lambda$.

En rassemblant :
\[
R=\mu\,(X-\lambda)^{2}\prod_{i=1}^{r}T_i,
\]
avec $\mu\neq 0$, les $T_i$ irréductibles deux à deux distincts et
$T_i(\lambda)\neq 0$ pour tout $i$ : c'est la décomposition
annoncée. (C'est exactement la forme exploitée aux questions 35 et
51--52 du sujet, avec $\K=\Fp$.)
\end{corrige}

\begin{pieges}
\begin{itemize}[leftmargin=1.4em, itemsep=2pt]
\item Le cas subtil : un facteur carré $P^{2}\mid T$ avec $\dg{P}=1$
ne permet pas d'appliquer la question 19 avec $\widetilde T=P$
(il faut $\dg{\widetilde T}\ge 2$) ; l'astuce consiste à
\emph{absorber} $(X-\lambda)$ en posant
$\widetilde T=(X-\lambda)P$, de degré $2$. Beaucoup de copies
oublient cette vérification d'hypothèse.
\item Oublier le scalaire $\mu$ (les $T_i$ sont choisis unitaires, le
coefficient dominant doit bien aller quelque part) ou oublier de
justifier $T_i(\lambda)\neq 0$.
\item Deux facteurs \og distincts \fg{} signifie ici non associés :
prendre les $T_i$ \emph{unitaires} rend la condition nette ;
l'éventuelle répétition $T_i=T_j$ ($i\neq j$) est déjà exclue par
l'argument $\alpha_i=1$.
\item Ne pas traiter (même en une phrase) le cas $T$ constant
($r=0$) : la formule reste vraie avec un produit vide.
\end{itemize}
\end{pieges}

% =====================================================================
\section*{Bilan : ce qu'il faut retenir de ces deux parties}
\addcontentsline{toc}{section}{Bilan}
% =====================================================================

\begin{conseil}
\begin{itemize}[leftmargin=1.4em, itemsep=2pt]
\item \textbf{Le fil conducteur.} La partie I fournit le dictionnaire
\og opérations élémentaires $=$ multiplications par des matrices de
$\GLn(\Z)$ \fg{} (questions 9--10), la structure de $\GLn(\Z)$
(question 6) et la description de $\On(\Z)$ (questions 7--8) : tout
cela ressert dans la forme normale de Smith (partie V) et dans la
conclusion (partie VII).
\item \textbf{La formule maîtresse de la partie II} est
$\dim\Ker(\varphi_{R,S})=\dg{R\wedge S}$ (question 15) : les
questions 16 à 20 en sont des corollaires directs. En particulier,
$\Delta(R)\neq 0\iff R\wedge R'=1$, et
$\dim\Ker(\psi_R)=1\iff R\wedge R'$ de degré $1$, ce qui caractérise
la forme $\mu(X-\lambda)^2\prod T_i$.
\item \textbf{Réflexe transversal :} en caractéristique quelconque,
bannir les arguments \og racines \fg{} et \og $\dg{R'}=\dg{R}-1$ \fg{}
au profit d'arguments de \emph{divisibilité} ; c'est le fil rouge des
pièges du sujet (VF4, 17, 18, puis partie IV).
\item \textbf{Gestion du temps.} Les quatre premières parties sont
indépendantes : le Vrai/Faux et la partie I sont des points
\og de cours \fg{} à sécuriser rapidement et proprement ; la partie
II demande davantage de soin (degrés, cas $S_1=0$, caractéristique)
et récompense les copies qui citent précisément la question 15.
\end{itemize}
\end{conseil}

\end{document}
